\section{Consistency and soundness conditions}
\label{sec:first-incompleteness-02}

\begin{definition}[Consistency]
A theory \(T\) is \emph{consistent} iff there is no sentence \(\sigma\) such
that both \(T\derives\sigma\) and \(T\derives\neg\sigma\).  In the fixed
classical explosive calculus, this is equivalent to
\(T\notderives\bottom\).  If \(T\derives 0\ne S0\), then it is also
equivalent to \(T\notderives 0=S0\); in particular this specialization applies
when \(T\supseteq\mathsf Q\).  Without the displayed derivability hypothesis,
nonderivability of \(0=S0\) is not a general test for consistency.
\end{definition}

Recall from \cref{def:theory-completeness} that \(T\) is deductively complete
when, for every sentence \(\sigma\) in its language, either
\(T\derives\sigma\) or \(T\derives\neg\sigma\).  This property of a theory is
distinct from semantic completeness of a calculus, while soundness and
completeness of the fixed calculus identify it with the quantified semantic
condition stated there.

\begin{definition}[\(\Sigma^0_1\)-soundness]
A theory \(T\) is \emph{\(\Sigma^0_1\)-sound} iff, whenever \(\sigma\) is a
\(\Sigma^0_1\) sentence and \(T\derives\sigma\), then
\(\mathcal N\entails\sigma\).
\end{definition}

\begin{definition}[Omega-consistency]
A theory \(T\) is \emph{omega-consistent} iff there is no formula
\(\varphi\) with \(\FV(\varphi)\subseteq\{x\}\) such that
\[
  T\derives\exists x\,\varphi(x)
  \quad\text{and}\quad
  T\derives\neg\varphi(\bar n)
  \text{ for every }n\in\N.
\]
\end{definition}

Both omega-consistency and \(\Sigma^0_1\)-soundness imply consistency, but
consistency alone is weaker.  Rosser's theorem will remove the stronger
hypothesis from the incompleteness conclusion.

